\clearpage
\VerdierPhysicalPageAnchor{0079}%
\VerdierSemanticAnchor{ch1-2-5-7-morphisms}
\noindent\hspace*{2em}--\ $X$ と $Y$ を $K(C)$ の二つの対象、
$f,g\colon X\rightrightarrows Y$ を $\operatorname{comp}(C)$ の二つの射とする。
「$f$ と $g$ はホモトピックである」という関係は
$\operatorname{Hom}_{\operatorname{comp}(C)}(X,Y)$ 上の同値関係であり、
$\operatorname{Hom}_{K(C)}(X,Y)$ を
$\operatorname{Hom}_{\operatorname{comp}(C)}(X,Y)$ のこの同値関係による
商として定める。

\VerdierSemanticAnchor{ch1-2-5-7-composition}
\noindent\hspace*{2em}--\ (2.5.5) から直ちに、
$\operatorname{comp}(C)$ の射の合成則は商に移り、
$K(C)$ の射の合成則を定めることが分かる。

\VerdierSemanticAnchor{ch1-2-5-7-additive-category}
(2.5.5) から直ちに、以上により確かに圏が定まり、
しかもこの圏が加法圏であることが分かる。

\VerdierSemanticAnchor{ch1-2-5-7-quotient-functor}
標準関手 $\operatorname{comp}(C)\xrightarrow{Q}K(C)$ は対象上恒等であり、
$\operatorname{comp}(C)$ の各射 $f$ にその同値類 $\dot f$ を対応させる。
この関手は加法的である。

\VerdierSemanticAnchor{ch1-2-5-7-homotopy-class}
$\operatorname{comp}(C)$ の射 $f$ の標準関手による像である
$K(C)$ の射 $\dot f$ を、$f$ の\emph{ホモトピー類}と呼ぶ。

\VerdierSemanticAnchor{ch1-2-5-7-name}
圏 $K(C)$ を $C$ の\emph{複体のホモトピー圏}と呼ぶ。

\VerdierSemanticAnchor{ch1-2-5-8}
\noindent\textbf{2.5.8.}\quad
標準関手 $\operatorname{comp}(C)\xrightarrow{Q}K(C)$ は次の性質をもつ：
$K(C)$ から圏 $A$ への任意の関手は、$Q$ との合成により
$\operatorname{comp}(C)$ から $A$ への関手を与える。したがって、
任意の圏 $A$ に対して関手

\VerdierSemanticAnchor{ch1-2-5-8-hom-functor}
\begin{equation*}
  \mathcal H\!om(Q,A)\colon
  \mathcal H\!om(K(C),A)
  \longrightarrow
  \mathcal H\!om(\operatorname{comp}(C),A)
\end{equation*}
を得る（$\mathcal H\!om$ は関手圏を表す）。

\VerdierSemanticAnchor{ch1-2-5-8-fully-faithful}
この関手は充満忠実であり、対象上単射である。その像は次の圏の充満部分圏である：

\VerdierSemanticAnchor{ch1-2-5-8-essential-image}
\begin{equation*}
  \mathcal H\!om(\operatorname{comp}(C),A).
\end{equation*}
この充満部分圏は、$\operatorname{comp}(C)$ のホモトピックな射の各対を
等しい射へ送る関手 $\operatorname{comp}(C)\longrightarrow A$ によって定まる。

\VerdierSemanticAnchor{ch1-2-5-8-uniqueness}
標準関手 $\operatorname{comp}(C)\xrightarrow{Q}K(C)$ を備えた圏 $K(C)$ は、
この性質によって一意な同型を除いて決まる。

\VerdierSemanticAnchor{ch1-2-5-9}
\noindent\textbf{2.5.9.}\quad
(2.5.5) と (2.5.8) から直ちに、$K(C)$ 上に唯一のシフト関手が存在し、
標準関手

\VerdierSemanticAnchor{ch1-2-5-9-translation}
\begin{equation*}
  \operatorname{comp}(C)\longrightarrow K(C)
\end{equation*}
はシフトと厳密に可換であることが分かる。

\VerdierSemanticAnchor{ch1-2-5-9-strict-z-category}
このシフト関手は $K(C)$ の自己同型であり、したがって $K(C)$ 上に
厳密な $\mathbb Z$-圏構造を定める。

\VerdierSemanticAnchor{ch1-2-5-9-strict-z-functor}
標準関手 $\operatorname{comp}(C)\longrightarrow K(C)$ は厳密な
$\mathbb Z$-関手である。
